\documentclass[10pt,a4paper]{amsart}
\usepackage[margin=19mm]{geometry}
\usepackage{amsmath,amssymb,mathtools}
\usepackage[scr=rsfs,cal=euler]{mathalfa}
\usepackage{xfrac}
\usepackage[unicode,hidelinks,bookmarks=false]{hyperref}
\newcommand{\ie}{i.e.\ }
\newcommand{\cf}{cf.\ }
\newcommand{\resp}{respectively\ }
\newcommand{\Subsection}{\subsection}
\providecommand{\nobreakdash}{\nobreak\hbox{-}}
\setlength{\emergencystretch}{2em}
\multlinegap0pt
\usepackage{bm}
\usepackage{bbm}
\usepackage{dsfont}
\usepackage{xspace}
\usepackage{cancel}
\usepackage{mathtools}
\usepackage{amsthm}
\makeatletter
\@ifundefined{thm}{\newtheorem{thm}{Thm}}{}
\@ifundefined{lemma}{\newtheorem{lemma}[thm]{Lemma}}{}
\makeatother
\def\P{\mathbb{P}}
\title[Optimal factor matchings for point processes on non-amenable]{Optimal factor matchings for point processes on non-amenable unimodular graphs}
\author{Yinon Spinka, Oren Yakir}
\date{Source: arXiv:2601.08983v1 (CC BY 4.0)}
\begin{document}
\maketitle

\noindent\emph{Source and licence.} Optimal factor matchings for point processes on non-amenable unimodular graphs. Version arXiv:2601.08983v1 (CC BY 4.0), distributed under the Creative Commons Attribution 4.0 International licence, as recorded in the arXiv licence metadata for this exact version and shown on the abstract page https://arxiv.org/abs/2601.08983v1. Full rights notice: licenses/arxiv-2601.08983-CC-BY-4.0.txt.\par\smallskip

\noindent\textit{Editorial scope.} Counted unit: the Lemma whose source label is \texttt{lem:distinct-point-counts} in the article (source0.tex lines 686--693, source label \texttt{lem:distinct-point-counts}), delivered together with the article's own complete original proof of it (source0.tex lines 718--760, one \texttt{proof} environment of 43 lines). No mathematics is written, repaired or re-derived: statement and proof are the article's own text line for line. The proof is detached from the statement in the source: the article states the result at lines 686--693 and prints its proof at lines 718--760, and the proof environment's own header names the statement, so the association is the article's own. Required context retained uncounted: none. Every definition, display and label of those lines is retained unchanged. Omitted from this package and not redistributed: 1--685, 694--717, 761--780. Nothing inside a retained excerpt is omitted or altered: every retained body is delivered exactly as the article's lines are written, so no omission receipt of any kind accompanies this package. Citations: the retained excerpts carry no citation command at all, so no bibliography entry of the article is transcribed and no reference is redistributed. Reference adaptation: none was needed and none was made; every reference inside the delivered unit resolves inside it, so no unresolved reference or citation is produced. Printed slips kept literal: no printed word, symbol, hypothesis or number of the counted statement or of its proof was altered. Layout and class: the article's own class is \texttt{article} with packages including amsmath, amsfonts, amssymb, amsthm, graphics, enumerate, mathtools, thmtools, hyperref, dsfont, graphicx, xcolor, bbm, fullpage; the wrapper is the lane's pinned \texttt{amsart} editorial wrapper at 10pt on A4, which restates the theorem-like environments the retained text needs and adds the article's own macro definitions that the retained text needs (\texttt{\textbackslash P}), with extra wrapper packages (bm, bbm, dsfont, xspace, cancel, mathtools). The delivered numbering is the unit's own. Assets: no retained span contains a figure environment, a graphics-inclusion command, a PDF-page inclusion, a table or a TikZ picture; no image file is copied into this package, the package declares no asset dependency and no image file is redistributed with it. Licence: version arXiv:2601.08983v1 of this article is distributed under the Creative Commons Attribution 4.0 International licence, as recorded in the arXiv licence metadata for this exact version and shown on the abstract page https://arxiv.org/abs/2601.08983v1. A full rights notice is delivered at licenses/arxiv-2601.08983-CC-BY-4.0.txt. Characters: every retained excerpt is pure ASCII, so the delivered engine, XeLaTeX with its default Latin Modern text font, reports no missing character.
\begin{lemma}\label{lem:distinct-point-counts}
    Let $G=(V,E)$ be a locally finite transitive connected graph such that $S_r(v) \neq S_r(u)$ for every distinct $u,v\in V$ and $r\ge 0$.    
    Let $\Pi$ be either a Poisson process, or a perturbed vertex set satisfying that there exists $c>0$ such that for every distinct $u,v\in V$ and $r\ge 0$ there exists $w \in V$ such that 
    \begin{equation} \label{eq:lem:distinct_point_count_condition_on_w}
        \text{\normalfont Var}\Big(\mathbf{1}_{\{X_{w} \in S_r(u) \setminus S_r(v)\}} \Big) \ge c.
    \end{equation}
    Then a.s.\ for any two distinct vertices $u,v\in V$ there exists $r\ge 0$ such that $|\Pi \cap S_r(v)| \neq |\Pi \cap S_r(u)|$.
\end{lemma}

\begin{proof}[Proof of Lemma~\ref{lem:distinct-point-counts}]
    We handle the case when $\Pi$ is a perturbed vertex set; the case when $\Pi$ is a Poisson process is similar (and much simpler).    
    Let $\{r_k\}$ be a subsequence that we choose later and denote by 
    \[
    E_k = \big\{ |\Pi\cap S_{r_k}(u)| \neq |\Pi\cap S_{r_k}(v)| \big\} \, .
    \]
    We want to show that $ \P(\bigcup_{k} E_k ) = 1$, which will follow immediately once we show that 
    \begin{equation} \label{eq:proof_of_lemma_distinct_points_lower_bound_on_conditional_prob}
        \inf_{k\ge 1} \P\big( E_k \mid E_1^c\cap \ldots \cap E_{k-1}^c \big) \ge \tfrac14 c \, .
    \end{equation}
    Fix some $k\ge 1$ and denote by $F_k = E_1^c\cap \ldots \cap E_{k-1}^c$. We can choose $r_k \ge 4 \max\{ r_{k-1}, \text{dist}(u,v)\}$ large enough so that 
    \begin{equation*}
        \P \big(\text{dist}(X_w,w) \ge \tfrac12 r_k \big) \le \tfrac12 \min\{c,\P(F_k)\} \, ,
    \end{equation*}
    for all $w$. By our assumption, there exists $w$ such that~\eqref{eq:lem:distinct_point_count_condition_on_w} holds with $r=r_k$.
    In particular, $\P(X_w \in S_{r_k}(u) \setminus S_{r_k}(v)) \ge c$, and hence, $\text{dist}(w, S_{r_k}(u) \setminus S_{r_k}(v)) \le r_k/2$. Thus, $\text{dist}(w,u) \ge r_k/2 \ge r_{k-1}$, and $\text{dist}(w,v) \ge \text{dist}(w,u) - \text{dist}(u,v) \ge r_k/4 \ge r_{k-1}$.
    Therefore,
    \begin{equation} \label{eq:proof_of_lemma_distinct_point_choice_of_r_k}
        \P \big(X_w \in B_{r_{k-1}}(u) \cup B_{r_{k-1}}(v) \big) \le \tfrac12 \min\{c,\P(F_k)\} \, ,
    \end{equation}
    
    Denoting by $\Pi^w =\{ X_v \, : \, v\not=w \}$, we have
    \begin{multline*}
    \mathbf{1}_{E_k} \ge \mathbf{1}_{\{X_w \in S_{r_k}(u) \setminus S_{r_k}(v)\}} \cdot \mathbf{1}_{\{|\Pi^w \cap S_{r_k}(v)| \le |\Pi^w \cap S_{r_k}(u)| \}} \\ +  \mathbf{1}_{\{X_w \notin S_{r_k}(u) \setminus S_{r_k}(v)\}} \cdot \mathbf{1}_{\{|\Pi^w \cap S_{r_k}(v)| > |\Pi^w \cap S_{r_k}(u)|\}} \, .
    \end{multline*}
    Noting that $X_w$ is independent of $\Pi^w$, we can take the conditional expectation over $F_k$ and observe from the above that
    \begin{equation} \label{eq:proof_of_lemmma_distinct_points_lower_bound_after_case_decomposition}
    \P(E_k \mid F_k )\ge \min\Big\{ \P\big(X_w \in S_{r_k}(u) \setminus S_{r_k}(v) \mid F_k \big) , \, \P\big(X_w \not \in S_{r_k}(u) \setminus S_{r_k}(v) \mid F_k \big)  \Big\} \, .    
    \end{equation}
    We now lower bound each of the terms on the right-hand side.
    The value of $X_w$ is irrelevant for the occurrence of $F_k$ on the event that $X_w \notin B_{r_{k-1}}(u) \cup B_{r_{k-1}}(v)$, and hence, for any $A \subset V \setminus (B_{r_{k-1}}(u) \cup B_{r_{k-1}}(v))$,
    \begin{align*}
        \P\big(X_w \in A \mid F_k \big) &= \P\big(X_w \in A \big) \cdot \frac{\P\big(F_k \mid X_w \in A \big)}{\P(F_k)} \\
        &= \P\big(X_w \in A \big) \cdot \frac{\P\big(F_k \mid X_w \not \in B_{r_{k-1}}(u) \cup B_{r_{k-1}}(v) \big)}{\P(F_k)} \\
        & \ge  \P\big(X_w \in A \big) \cdot \frac{\P(F_k) - \P\big( X_w \in B_{r_{k-1}}(u) \cup B_{r_{k-1}}(v) \big)}{\P(F_k)} \stackrel{\eqref{eq:proof_of_lemma_distinct_point_choice_of_r_k}}\ge \frac12 \, \P\big(X_w \in A \big) \, .
    \end{align*}
    On the other hand, taking $A=S_{r_k}(u) \setminus S_{r_k}(v)$, we get that $\P\big(X_w \in S_{r_k}(u) \setminus S_{r_k}(v) \mid F_k \big) \ge c/2$.
    Taking $A=(S_{r_k}(u) \setminus S_{r_k}(v))^c \setminus (B_{r_{k-1}}(u) \cup B_{r_{k-1}}(v))$, we get that
    \begin{align*}
       \P\big(X_w \notin S_{r_k}(u) \setminus S_{r_k}(v) \mid F_k \big) &\ge \P\big(X_w \in  (S_{r_k}(u) \setminus S_{r_k}(v))^c \setminus (B_{r_{k-1}}(u) \cup B_{r_{k-1}}(v))  \mid F_k \big) \\ & \ge \frac{1}{2} \P\big(X_w \in  (S_{r_k}(u) \setminus S_{r_k}(v))^c \setminus (B_{r_{k-1}}(u) \cup B_{r_{k-1}}(v))  \big) \\ &\ge \frac{1}{2} \Big(\P\big(X_w \in  (S_{r_k}(u) \setminus S_{r_k}(v))^c   \big) -  \P \big(X_w\in (B_{r_{k-1}}(u) \cup B_{r_{k-1}}(v) \big) \Big) \\  &\stackrel{\eqref{eq:proof_of_lemma_distinct_point_choice_of_r_k}}\ge \frac14 c . 
    \end{align*}
    Thus, \eqref{eq:proof_of_lemma_distinct_points_lower_bound_on_conditional_prob} follows from~\eqref{eq:proof_of_lemmma_distinct_points_lower_bound_after_case_decomposition}, and we are done. 
\end{proof}

\end{document}
